\documentclass[../main.tex]{subfiles}
\begin{document}

\section{Introducción}

% Objeto
\subsection{Objeto}
Analizar las posibles alternativas al propano como fluido de trabajo para ciclos de bombas de calor simples, de forma teórica y experimental, que incrementen la eficiencia energética sin comprometer excesivamente otros parámetros importantes.

% Alcance
\subsection{Alcance}
\#\#\# Fer-ho quan haja acabat el experimental \#\#\#

% Antecedentes y justificación
\subsection{Antecedentes}
% Sistemas convencionales de generación de calor
En la actualidad, el 62\% del calor destinado a la calefacción de edificios se genera a partir de combustibles fósiles, el 26\% procede de biomasa y solamente el 12\% proviene de renovables modernas como bombas de calor, energía termosolar o geotermia. En el caso de la calefacción urbana, esta dependencia es aún mayor, ya que alrededor del 90\% de las instalaciones siguen basándose en combustibles fósiles (IRENA 2023) \cite{IRENA2023}.

% BOMBAS DE CALOR FRENTE A CALDERAS
No obstante, diversos estudios señalan que la sustitución de los sistemas de calefacción basados en combustibles fósiles por sistemas con bomba de calor permitiría reducir de forma significativa tanto el consumo energético como el impacto ambiental (Mayer 2024) \cite{Meyer_2024}.

% FUNCIONAMIENTO BOMBA DE CALOR
Las bombas de calor por compresión se basan en la transferencia de calor de un lugar a otro. El ciclo hace circular un fluido de trabajo por su interior y, de forma cíclica y cerrada, se evapora, comprime, condensa y expande absorbiendo y liberando calor en el acto.

\begin{figure}[ht]
    \centering
    \includegraphics[width=200pt]{imagenes/ciclo_basico.png}
    \caption{Ciclo básico de compresión de vapor}
    \label{fig:ciclo_basico}
\end{figure}

Como se puede observar en la \autoref{fig:ciclo_basico}, este ciclo tiene cuatro componentes principales: un evaporador, un compresor, un condensador y un dispositivo de expansión. De este último se puede aprovechar el calor cedido por el condensador o el calor absorbido por el evaporador, tanto para calentar como para enfriar, respectivamente, lo que lo hace muy útil tanto a nivel doméstico como a nivel industrial.

% PROBLEMAS BOMBAS DE CALOR
A día de hoy, solamente el sector de la calefacción en los edificios mediante bombas de calor presenta un 10\% de las emisiones de \begin{math} \left[\mathrm{CO_2}\right]_{\mathrm{eq}} \end{math} a la atmósfera (sin considerar otros gases de efecto invernadero). De estas emisiones, un tercio son emisiones directas, es decir, provienen de la liberación de fluidos de trabajo con alto potencial de calentamiento global a la atmósfera, mientras que los otros dos tercios están asociados a emisiones indirectas (IEA 2022) \cite{IEA2022FutureHeatPumps}.

% MARCO HISTÓRICO-REGULATIVO
\subsection{Marco histórico-regulativo}

El sector de las bombas de calor se ha desarrollado mucho desde su descubrimiento debido a grandes avances tecnológicos, termodinámicos, químicos y medioambientales. El primer sistema de bomba de calor por compresión fue descrito por Oliver Evans en el 1805, aunque no fue hasta el 1830 que John Hague lo construyó usando etil éter como fluido de trabajo.

Es, a partir de entonces, cuando consideramos que empieza la historia de las bombas de calor. Diferenciamos 4 principales generaciones, marcadas por el fluido de trabajo usado en cada una de ellas.

\subsubsection{Primera generación}
La primera generación (1830-1930) se caracterizó por centrarse en lograr el funcionamiento de las bombas de calor más que en su eficiencia u optimización, debido al limitado conocimiento técnico de la época. Por ello, los fluidos de trabajo empleados eran principalmente sustancias disponibles de forma natural o fácilmente accesibles, como el amoníaco (NH\tsub{3}), el dióxido de carbono (CO\tsub{2}), el propano (CH\tsub{3}-CH\tsub{2}-CH\tsub{3}), el dióxido de azufre (SO\tsub{2}) y el dimetil éter (CH\tsub{3}O-CH\tsub{3}).


La seguridad pasó a un segundo plano y, en consecuencia, ocurrieron varios accidentes laborales que hicieron replantearse las medidas mínimas de seguridad en la industria de las bombas de calor.



\subsubsection{Segunda generación}
La segunda generación surgió con el trabajo de Thomas Midgley Jr. y Albert L. Henne sobre los halogenúros orgánicos volátiles. Los investigadores estaban buscando la solución a los problemas de los fluidos de trabajo anteriores cuando encontraron la familia de los clorofluorocarburos (CFC) y, añadiendo hidrógeno a la mezcla obtuvieron los hidroclorofluorocarburos (HCFC).

Estas moléculas consiguieron reducir la inflamabilidad, toxicidad y aumentar la estabilidad y compatibilidad con los materiales presentes en las bombas de calor. Algunos fluidos de trabajo que dominaron esta etapa fueron el R11 (CCl\tsub{3}F), el R12 (CCl\tsub{2}F\tsub{2}), el R22 (CHClF\tsub{2}) y el R502 (mezcla de HCFC y CFC).

Sin embargo, tras una época de dominación de los CFC y HCFC, los investigadores Rowland y Molina, descubrieron que si estas moléculas llegaban suficientemente alto en la atmósfera, en contacto con la radiación ultravioleta, el cloro se desprendía y empezaba una reacción en cadena por catálisis, donde un solo átomo de cloro podía destruir miles de moléculas de ozono.

Para resolver la problemática de la desaparición de la capa de ozono, en el año 1987 se decidió firmar el Protocolo Montreal, un acuerdo internacional que imponía una reducción progresiva de los fluidos de trabajo que afectaban a la capa de ozono hasta llegar a su eliminación.


\subsubsection{Tercera generación}

Debido a la prohibición de los CFC y HCFC, esta generación empezó con el desarrollo de los hidrofluorocarburos (HFC), que no contienen moléculas de cloro, entre ellos destacan el R134a (CF\tsub{3}-CH\tsub{2}F), el R404A (mezcla de HFC) y el R507A (mezcla de HFC). Estos fluidos de trabajo, aunque no destruyen la capa de ozono, tienen otro efecto nocivo sobre el medioambiente, ya que, al contener moléculas de flúor, la absorción de radiación infrarroja se ve favorecida agravando el efecto invernadero.

Posteriormente, el Protocolo Kioto de 1997 establecería un límite para estos gases con alto potencial de calentamiento global (GWP, por sus siglas en inglés) y asentaría una base para las futuras normativas respecto a los gases de efecto invernadero.



\subsubsection{Cuarta generación (actual)}

En 2005 entraría en vigor el protocolo Kioto y una década después, a nivel europeo se establecería la ley F-gas (UE 517/2014) \cite{EU2014_517}, que reforzaría drásticamente las restricciones en el uso de HFC como se puede ver en la \autoref{fig:max_CO2_f-gas}.

Esta generación continuaría con la aparición de las hidrofluoroolefinas (HFO), unos gases fluorados pero con un doble enlace entre los carbonos, consecuentemente, su duración en la atmósfera es muy reducida y, por lo tanto, no afectan a la capa de ozono y tienen un GWP muy bajo.

Los principales HFO que aparecen son el R1234yf (CH\tsub{2}=CF-CF\tsub{3}), que sustituye al HFC R134a, y el R1234ze (CHF=CH-CF\tsub{3}), que sustituye tanto al R134a como al R410A (ambos HFC).

Además, la transición de HFC a HFO se vio potenciada por la actualización de la ley F-gas en el año 2024 (UE 573/2024) \cite{EU2024_573}, que sería más estricta con las cantidades máximas permitidas de gases con alto GWP en el mercado. Véase la \autoref{fig:max_CO2_f-gas}.

\begin{figure}[ht]
    \centering
    \includegraphics[width=0.7\textwidth]{imagenes/max_CO2_f-gas_v2.png}
    \caption{Máxima cantidad de CO\tsub{2} eq de HFC entrante en el mercado según F-gas}
    \label{fig:max_CO2_f-gas}
\end{figure}

Sin embargo, en estos últimos años, los HFO han sido fuente de preocupación por su capacidad de emitir sustancias perfluoroalquiladas y polifluoralquiladas (PFAS) y estas, aunque no impacten negativamente al medioambiente, conllevan un riesgo para la salud humana y animal.

Es por este motivo que ha crecido el interés en los gases naturales, ya que, aunque no son los mejores en cuanto a rendimiento energético, se sabe que son los que menos impacto tendrán tanto en la salud como a nivel medioambiental.


% Motivación y objetivos
\subsection{Motivación y objetivos}

% Estructura
\subsection{Estructura}


\newpage
\end{document}